\section{The single-clock compatibility obstruction}

The modeled clock used modulus $4LQ$, cutoff $\lfloor Lq/Q\rfloor$, and profile
argument $mQ/(Lq)$.  The next theorem characterizes exact transfer without an
asymptotic approximation.

\begin{theorem}[Compatibility if and only if]\label{thm:compatibility}
For a nonzero depth $L$, the physical row \eqref{eq:physical-row} and the modeled row agree in modulus and
in cutoff/profile scale for every multiplier and shell prime if and only if
\begin{equation}
 h=4LQ,\qquad H=4Q^2.                                  \label{eq:compatibility}
\end{equation}
\end{theorem}

\begin{proof}
Equality of moduli gives $h=4LQ$.  Equality of the profile arguments for a nonzero
multiplier gives $H/h=Q/L$.  Substitution yields $H=4Q^2$.  Conversely, the two
identities in \eqref{eq:compatibility} make the modulus, cutoff, profile argument, and residue condition
identical.
\end{proof}

At the V59 scales,
\begin{equation}
 \frac{4Q^2}{H}=4x^{2/3-21/32}=4x^{1/96}\longrightarrow\infty. \label{eq:clock-gap}
\end{equation}
Thus exact attachment to the single-clock family is refuted in this scope.  Matching
the physical depth makes the modeled modulus too large by \eqref{eq:clock-gap}; matching the modulus
makes its multiplier depth too small by the same factor.  This is not a constant-loss
bookkeeping discrepancy: the mismatch grows on precisely the conductor-gap scale
$Q^2/H$.

Theorem~\ref{thm:compatibility} does not invalidate the modeled clock as an abstract
collision problem.  It invalidates only the automatic identification of one such
clock with one physical V59 row.
